\documentclass[12pt]{article}

\usepackage[utf8]{inputenc}
\usepackage[T1]{fontenc}
\usepackage{amsmath}
\usepackage{amssymb}
\usepackage{graphicx}
\usepackage{booktabs}
\usepackage{hyperref}
\usepackage[round]{natbib}
\usepackage[margin=1in]{geometry}

\title{The Accumulation--Renewal Dilemma: Why Homogeneous Knowledge Degrades Intelligence and Only Asymmetric Renewal Restores It}
\author{Jiyeon Yoon\\
Independent Researcher\\
\texttt{heznpc@gmail.com}}
\date{}

\begin{document}

\maketitle

\begin{abstract}
Intelligent systems improve by accumulating knowledge, yet accumulation also
induces \emph{entrenchment}: a growing insensitivity to evidence that conflicts
with what has already been stored. We argue that this accumulation--renewal
dilemma recurs across substrates that have no mechanism in common---human
confirmation bias, paradigm lock-in in science \citep{kuhn1962structure},
the incumbent's trap in firms \citep{christensen1997innovator}, and context
entrenchment in large language models (LLMs)---and that each independently
arrives at the same structural resolution: \emph{asymmetric renewal}, the
controlled injection of a perspective that did not accumulate the same context.
We make this a design principle---the \textbf{Principle of Least Context}: give
an agent the minimum context its task requires, because additional context is a
source of bias before it is a source of capability. We position the LLM context
window as the first cheaply controllable testbed for the dilemma, and report a
small pilot suggesting that \emph{context} asymmetry (deep vs.\ fresh sessions)
can outperform \emph{model} diversity (cross-model debate). We are deliberate
about the pilot's weakness: with three tasks and a single run per condition, an
apparent ``direction-of-asymmetry'' effect is confounded with deep-session model
competence and cannot be claimed; we specify a pre-registered, confound-
controlled experiment that would adjudicate it. The contribution of this paper
is the framework and the principle, not the pilot. Code, task checklists, and
the experimental design are released for independent replication.
\end{abstract}

\section{Introduction}

Programmers know a particular experience: after hours stuck on a bug, you write
down what you have learned and \emph{start over in a clean session}. Often the
fresh start succeeds where persistence failed---the accumulated trail of
attempts had itself constrained the search. The same pattern recurs at larger
scales. Experienced incumbents stall in their own conventions while a newcomer,
unburdened by the accumulated context, reshapes the field
\citep{christensen1997innovator}. \emph{Accumulation is itself a source of bias,
and renewal tends to arrive from outside.}

Human progress is built on accumulated knowledge, yet the paradox that
accumulated knowledge impedes its own revision is repeatedly observed. Paradigm
shifts proceed less by persuading entrenched scientists than by generational
turnover \citep{kuhn1962structure}; quantitatively, the premature death of an
eminent scientist is followed by a measurable rise in non-collaborator entry and
new directions in the affected field \citep{azoulay2019}. We call the general
phenomenon the \emph{accumulation--renewal dilemma}: accumulating knowledge
raises capability while inducing entrenchment that lowers responsiveness to new
evidence.

\paragraph{Stochastic branching.} An LLM answers the same prompt differently
across sessions. A user who stays in one session lets the stochastic sample of
the first response constrain the entire downstream trajectory. Spawning $N$
agents that share one accumulated context does not fix this: they sample $N$
times from the \emph{same} biased distribution. Recent results make the limit
precise---under symmetric information, debate among agents induces a martingale
over belief trajectories, so debate alone does not improve expected correctness
\citep{choi2025debate}. Breaking the martingale requires an \emph{asymmetry}
between the participants \citep{liu2026acemad,young2026divergence}.

\paragraph{The LLM as a controllable analog.} Where human generational turnover
takes decades, an LLM session can be renewed instantly and the conflict between
old and new context can be staged under control. This makes the context window
the first setting in which the accumulation--renewal dilemma is cheaply
manipulable rather than merely observed.

\paragraph{Contributions.}
(i) We formalize the accumulation--renewal dilemma and the Principle of Least
Context (\S\ref{sec:framework}).
(ii) We report a pilot indicating that context asymmetry can outperform model
diversity, and we are explicit about what the pilot cannot show
(\S\ref{sec:pilot}).
(iii) We specify a pre-registered experiment that separates ``direction of
asymmetry'' from deep-session competence (\S\ref{sec:pilot}).
The framework and principle are the contribution; the empirical material is a
preliminary probe.

\section{Related Work}

\paragraph{Context degradation.} Performance degrades with context length even
under perfect retrieval \citep{du2025context,liu2023lost}, and effective context
capacity is well below advertised windows \citep{chroma2025contextrot,
scaling2026paradox}. Multi-turn accumulation reshapes model confidence
\citep{harshavardhan2026sacd} and causes models to ``get lost''
\citep{laban2025lost}; interaction context tends to increase sycophancy
\citep{jain2025sycophancy}; prompt-based debiasing of anchoring is
statistically ineffective \citep{feng2026anchoring}.

\paragraph{Multi-agent debate.} Symmetric debate is a martingale
\citep{choi2025debate}; asymmetric cognitive potential can restore positive
drift \citep{liu2026acemad}, and debate's value scales with knowledge divergence
between participants \citep{young2026divergence}. Cross-context review in a fresh
session improves output quality \citep{song2026ccr}, while adding rounds adds
noise \citep{song2026dccr}. Related lines reframe debate cooperatively
\citep{chen2025colmad}, with typed epistemic acts
\citep{debate2026deliberation}, with self-reflective context adjudication
\citep{srdcr2025}, or study when debate genuinely helps
\citep{wu2025debate,m2cl2026}. Entrenchment in debate itself is documented
\citep{oh2025dreamad}.

\paragraph{Scaling without renewal.} Large agent simulations scale population
size \citep{pang2025agentsociety}; collective bias can emerge even when no
individual agent is biased \citep{ashery2025conventions}; behavior degrades over
extended interaction \citep{rath2026agentdrift}. At the training scale,
recursively training on synthetic output collapses quality and diversity
\citep{shumailov2024collapse,alemohammad2024mad}.

\paragraph{Cognitive and biological background.} The dilemma's vocabulary is
borrowed, not claimed as mechanism: anchoring-and-adjustment
\citep{epley2006anchoring}, retrieval-induced and desirable forgetting
\citep{storm2012forgetting,slamecka1978generation,roediger2006test}, incubation
\citep{sio2009incubation}, reconstructive memory \citep{bartlett1932remembering},
sleep-dependent synaptic down-selection \citep{tononi2014plasticity}, energetic
limits that force selective plasticity
\citep{tyree2023competitive,li2024energyefficiency}, and forgetting as resource
management \citep{liu2024forgettingsurvey,xie2026sleepgate}.

\section{Framework: The Accumulation--Renewal Dilemma}
\label{sec:framework}

\paragraph{Definitions.} Let an agent $a$ accumulate context over a session.

\begin{description}
\item[Definition 1 (Context Lifespan).] $L(a)$ is the total accumulated context
(in tokens) from session start to the present.
\item[Definition 2 (Context Entrenchment).] $E(a,L)$ is the probability that
$a$, at lifespan $L$, produces a response consistent with its own initial
stochastic sample.
\end{description}

\begin{description}
\item[Hypothesis 1 (Entrenchment Monotonicity).] $E(a,L)$ is non-decreasing
in $L$.
\item[Hypothesis 2 (Asymmetric Advantage).] A structured debate between two
agents with \emph{different} $L$ improves decision quality over a debate between
two agents with the \emph{same} $L$.
\item[Hypothesis 3 (Generational Innovation).] Periodically introducing a fresh
agent into a structured conflict yields better long-run decisions than a
homogeneous population.
\end{description}

\paragraph{Where the analogy holds and breaks.} The analogy is structural, not
mechanistic, and we mark its limits explicitly. It \emph{holds} for
anchoring, consistency-seeking, and the ``fresh eyes'' advantage. It
\emph{breaks} in that an LLM's ``fresh'' session is literally zero-context
(a human newcomer brings a different lifetime of experience), session renewal
is non-destructive (human turnover is irreversible), and human institutions add
political dynamics absent from the LLM case. Convergence across substrates is
suggestive of a shared \emph{design constraint}, not evidence of a shared
mechanism, and the set of substrates is curated rather than exhaustive.

\paragraph{The Principle of Least Context.} The security principle of least
privilege \citep{saltzer1975protection} has cognitive force here.

\begin{quote}
\textbf{Principle of Least Context.} Provide an agent with the minimum context
its task requires. Additional context is a source of bias before it is a source
of capability.
\end{quote}

This runs against the industry assumption that a longer context window is
strictly better.

\paragraph{Low species diversity.} A complementary reason renewal must be
engineered: LLMs share an architecture \citep{vaswani2017attention} and overlap
heavily in training data, and their representations appear to converge
\citep{huh2024platonic}. ``Model diversity'' is therefore shallow; differences
in \emph{context} may be a more reliable source of useful disagreement than
differences in model.

\section{A Pilot Probe, and the Confound It Cannot Resolve}
\label{sec:pilot}

\paragraph{Setup.} As a first probe we ran three long-context tasks (each
seeded with a plausible-but-wrong early assumption) under six conditions using a
stronger and a weaker model, with a structured debate
(\textsc{independent}$\rightarrow$\textsc{position}$\rightarrow$\textsc{challenge}
$\rightarrow$\textsc{convergence}). Recall is the fraction of a fixed checklist
of issues surfaced. Results are in Table~\ref{tab:pilot}.

\begin{table}[h]
\centering
\begin{tabular}{lcc}
\toprule
Condition & Diversity type & Recall \\
\midrule
Single (strong) & baseline & 83.3\% \\
Single (weak) & baseline & 58.9\% \\
Cross-model symmetric & model & 74.4\% \\
Same-model asymmetric (weak) & context & 90.0\% \\
Cross-model asymmetric (strong deep $\rightarrow$ weak fresh) & model+context & 93.9\% \\
Cross-model asymmetric (weak deep $\rightarrow$ strong fresh) & model+context & 60.0\% \\
\bottomrule
\end{tabular}
\caption{Pilot results ($n=3$ tasks, single run per condition). Numbers are
reported to one decimal place only to match the underlying counts; they are not
precise estimates---see the caveat below.}
\label{tab:pilot}
\end{table}

\paragraph{What the pilot suggests.} Two patterns are consistent with the
framework. First, model diversity alone did not help (cross-model symmetric
$74.4\% < 83.3\%$ single-strong), echoing the symmetric-debate martingale
\citep{choi2025debate}. Second, same-model context asymmetry did help the weaker
model ($58.9\% \rightarrow 90.0\%$), consistent with asymmetry as the active
ingredient \citep{liu2026acemad,song2026ccr}.

\paragraph{What the pilot cannot show.} The eye-catching contrast---strong-deep
$\rightarrow$ weak-fresh ($93.9\%$) vs.\ weak-deep $\rightarrow$ strong-fresh
($60.0\%$)---is \textbf{confounded}: the ``direction'' of asymmetry covaries with
which model held the deep role, and the deep model's single-session baselines
($83.3\%$, $58.9\%$) already track the debate outcomes. A competing explanation
needs no notion of direction at all: \emph{final quality is dominated by the deep
session's competence.} With $n=3$ and one run per cell, this null cannot be
rejected, and we therefore make \emph{no} directional claim. We treat
``direction matters'' as an open question, not a finding.

\paragraph{A test that would resolve it.} We pre-register a
$2\times2$ factorial---deep $\in\{\text{strong},\text{weak}\}\times$ fresh
$\in\{\text{strong},\text{weak}\}$---with same-model asymmetry cells (which
remove model identity entirely) and compute-control baselines (two strong agents
sharing context; one strong model sampled twice). The decisive quantity is the
\emph{deep$\times$fresh interaction} in
$\text{recall}\sim\text{deepStrong}*\text{freshStrong}+(1\,|\,\text{task})$ over
$\geq 30$ tasks with $\geq 3$ repeats and a blind third-model judge. If the
interaction is null, the directional reading is rejected as a deep-competence
artifact; if it survives \emph{and} exceeds the compute controls, the direction
effect is real. Either outcome is informative. The full protocol, task
checklists, and analysis are released with this paper.

\section{Discussion}

\paragraph{The incumbent trap.} The dilemma reframes the innovator's dilemma
\citep{christensen1997innovator} mechanistically: an incumbent's ``experience''
is accumulated context, a newcomer's ``naivety'' is fresh context, and the
incumbent rationally misses disruption not from lack of ability but because
accumulated context constrains the search space. The same dynamic appears at the
session level: a deep session anchored on a long rationale fails to question it,
while a fresh session flags the obvious objection immediately.

\paragraph{Persistent memory as a confabulation amplifier.} The dilemma predicts
a concrete failure in deployed tooling. Cross-session ``memory'' that injects a
fixed store into every new session is exactly the homogeneous-accumulation
pattern: opening a fresh session no longer yields fresh context, and any
hallucination written to the store becomes the next session's trusted premise.
This is a within-system analog of model collapse
\citep{shumailov2024collapse}---a session-to-session echo of the chat-chamber
effect, in which users come to trust AI-generated falsehoods
\citep{jacob2025chatchamber}. The framework predicts the mitigations that help:
periodic expiry (an engineered Planck principle) and citation verification
(asymmetric checking). We describe the pattern at the level of mechanism; the
specific deployed-product incidents that motivated it require primary-source
confirmation before they can be cited as evidence.

\paragraph{Design implication.} ``More agents, longer context'' is a limited
direction. A few sessions with deliberately different context depth, surfaced
through a structured protocol, may decide better than a large homogeneous
population---a hypothesis the experiment above is designed to test.

\section{Limitations}

The completed empirical material is a single pilot ($n=3$ tasks, one run per
condition, two models, author-constructed checklists, a single judge); its
differences are within plausible run-to-run variance, and its headline contrast
is confounded as discussed in \S\ref{sec:pilot}. The cross-substrate framing is
an analogy whose mechanisms differ across substrates, and the substrate set is
curated. Larger validation requires community-scale execution; the released
design and harness are intended to enable it.

\section{Conclusion}

The accumulation--renewal dilemma names a pattern shared by minds, sciences,
firms, and language models: homogeneous accumulation entrenches, and only the
controlled injection of a non-accumulated perspective renews. The Principle of
Least Context turns this into a design rule. Whether the \emph{direction} of
asymmetry carries independent weight is left as a sharply specified, falsifiable
question rather than asserted from a pilot.

\bibliographystyle{plainnat}
\bibliography{references}

\end{document}
